\chapter{The Atmosphere: Multi-Level Spectral Dynamical Core}\label{chap:ch09_atmosphere_multilevel}

% --- local symbols (kept minimal; not in the shared preamble) -------
\providecommand{\sig}{\sigma}                 % vertical sigma coordinate
\providecommand{\sigd}{\dot{\sigma}}          % vertical sigma velocity
\providecommand{\dvg}{\delta}                 % horizontal divergence (spectral prognostic)
\providecommand{\Tref}{\overline{T}}          % reference temperature profile T_l(sigma)
\providecommand{\Tvar}{T'}                     % temperature variation about reference
\providecommand{\Teq}{T_{\mathrm{eq}}}        % radiative-equilibrium temperature
\providecommand{\qhum}{q}                      % specific humidity tracer
\providecommand{\Lv}{L_v}                      % latent heat of vaporization

The single-level barotropic atmosphere of Chapter~\ref{chap:ch10_atmosphere_singlelevel}
has a fatal structural limitation: with no vertical degree of freedom it cannot
support \emph{baroclinic instability}, the conversion of the equator--pole
temperature gradient into eddy kinetic energy that builds mid-latitude storm
tracks, the eddy-driven jet, and the surface westerlies. \model{}'s production
atmosphere therefore replaces it with a genuine multi-level spectral
primitive-equation dynamical core. Rather than re-implement a spectral transform
core from scratch, \model{} wraps Google's \dino{} dycore --- the differentiable
\jax{} spectral primitive-equation solver that underlies NeuralGCM
\citep{kochkov2024neuralgcm} --- and supplies the column physics
(thermal forcing, drag, and a moisture cycle) as additional explicit tendencies.
Everything in this chapter lives in
\codref{chronos\_esm/atmos/dino\_atmos.py}{the \code{DinoAtmosphere} wrapper class}.

This is a deliberate architectural choice. The dycore handles the parts that are
hard to get right and easy to get wrong --- the spherical-harmonic transforms,
the semi-implicit treatment of gravity waves, and tracer advection --- while
\model{} owns the physics that couples the atmosphere to the rest of the Earth
system. The result is differentiable end to end, so the adjoint of an atmospheric
trajectory can be taken with respect to forcing or initial conditions
(Chapter~\ref{chap:ch03_differentiable}).

\begin{remark}
The standalone single-file \code{spectral.py} in the same directory
(\codref{chronos\_esm/atmos/spectral.py}{FFT-in-$\lon$ / finite-difference-in-$\lat$
transforms, \code{solve\_helmholtz}}) is a \emph{pseudo}-spectral hybrid written
for the legacy single-level model and the QTCM experiments. Its header is candid:
a full pure-\jax{} Legendre transform is ``1000+ lines,'' so it uses an exact
zonal FFT but fourth-order finite differences in latitude. The multi-level
atmosphere described here does \emph{not} use it; \dino{} provides true
spherical-harmonic transforms. We mention \code{spectral.py} only to map the
repository honestly --- its \code{solve\_helmholtz} is the building block of the
semi-implicit scheme the legacy model never wired in.
\end{remark}

\section{Governing equations in $\sig$ coordinates}

The dycore integrates the hydrostatic primitive equations on the sphere in a
terrain-following vertical coordinate
\begin{equation}
  \sig \;=\; \frac{\pres}{\pres_s},
\end{equation}
where $\pres_s(\lon,\lat,t)$ is the surface pressure, so $\sig$ runs from $0$ at
the model top to $1$ at the surface and the lower boundary $\sig=1$ is a
coordinate surface even over orography \citep[Ch.~12]{vallis2017}. \model{} uses
\code{SigmaCoordinates.equidistant(24)}: $24$ equally spaced layers of thickness
$\Delta\sig = 1/24$, with layer centres $\sig_l$ stored in \code{self.sigma}
(\codref{chronos\_esm/atmos/dino\_atmos.py}{\code{DinoAtmosphere.\_\_init\_\_},
vertical grid}). The horizontal grid is the spectral T31 triangular truncation
(\code{Grid.T31}, $\approx 3.75^\circ$ Gaussian nodal mesh).

Following standard spectral practice, the horizontal momentum is carried not as
$(\uu,\vv)$ but as its rotational and divergent parts --- the relative vorticity
$\vort$ and the horizontal divergence $\dvg$:
\begin{equation}
  \vort \;=\; \kvec\cdot(\curl\,\uvec), \qquad
  \dvg  \;=\; \diver\,\uvec .
\end{equation}
The four prognostic fields, stored as complex spherical-harmonic coefficients in
the \dino{} \code{State} object, are
\begin{equation}
  \big(\;\vort,\;\dvg,\;\Tvar,\;\ln\pres_s\;\big),
\end{equation}
where $\Tvar = T - \Tref_l$ is the temperature \emph{variation} about a fixed
reference profile $\Tref_l$ (the layer-mean temperatures $\approx 288$~K returned
by \code{isothermal\_rest\_atmosphere}), and $\ln\pres_s$ is the log of surface
pressure (carrying $\ln\pres_s$ rather than $\pres_s$ linearises the surface-
pressure tendency and keeps $\pres_s>0$ exactly). A specific-humidity tracer
$\qhum$ is advected alongside them. The continuous evolution is, schematically,
\begin{align}
  \pd{\vort}{t} &= -\,\kvec\cdot\curl\!\big(\,\text{absolute-vorticity flux}\,\big)
                   + F_\vort, \\
  \pd{\dvg}{t}  &= \kvec\cdot\curl(\dots) - \lap\!\Big(\tfrac{1}{2}|\uvec|^2 + \Phi\Big)
                   + F_\dvg, \\
  \pd{\Tvar}{t} &= -\uvec\cdot\grad T - \sigd\,\pd{T}{\sig}
                   + \frac{\kappa T\,\omega}{\pres} + F_T, \\
  \pd{\ln\pres_s}{t} &= -\,\frac{1}{\pres_s}\!\int_0^1 \diver(\pres_s\,\uvec)\dd\sig ,
  \label{eq:lnps}
\end{align}
with geopotential $\Phi$ from hydrostatic balance, $\kappa = R_d/c_p$, vertical
$\sig$-velocity $\sigd$ diagnosed from the continuity constraint, and
$\omega = \Dt{\pres}$. The $F$ terms are the \model{}-supplied forcings of
Section~\ref{sec:ch09forcing}. The exact discrete tendencies are \dino{}'s
\code{PrimitiveEquations.explicit\_terms} and \code{implicit\_terms}; \model{}
constructs the operator with the reference profile, the modal orography, and the
coordinate/spec objects (\codref{chronos\_esm/atmos/dino\_atmos.py}{\code{pe.PrimitiveEquations(...)}}).
Diagnostic $(\uu,\vv)$ are recovered from $(\vort,\dvg)$ by inverting the
streamfunction/velocity-potential relations inside
\code{pe.compute\_diagnostic\_state}, which returns $\uu\cos\lat$ and $\vv\cos\lat$;
\model{} divides by $\cos\lat$ to get physical winds
(\codref{chronos\_esm/atmos/dino\_atmos.py}{\code{DinoAtmosphere.diagnostics}}).

\section{The spectral transform method}

The spectral method represents each horizontal field as a truncated sum of
spherical harmonics $Y_l^m(\lat,\lon) = P_l^m(\sin\lat)\,e^{im\lon}$,
\begin{equation}
  f(\lon,\lat) \;=\; \sum_{m=-M}^{M}\;\sum_{l=|m|}^{L} \hat f_l^m\, Y_l^m(\lat,\lon),
  \qquad \text{T31: } L = M = 31 .
\end{equation}
Two operators motivate the choice. First, horizontal derivatives are
\emph{exact and local} in spectral space: the Laplacian is diagonal,
\begin{equation}
  \lap Y_l^m \;=\; -\,\frac{l(l+1)}{\Erad^2}\,Y_l^m ,
  \label{eq:laplacian-eig}
\end{equation}
so the Poisson inversions needed to get $\uvec$ from $(\vort,\dvg)$ and $\Phi$
from temperature become mere multiplications --- no iterative solver, no polar
finite-difference singularity. Second, the nonlinear products (advection, kinetic
energy) are cheapest to evaluate pointwise on the grid. The dycore therefore
\emph{transforms} between the two representations every step: it computes
derivatives spectrally, transforms to the Gaussian nodal mesh
(\code{to\_nodal}), forms products there, and transforms back (\code{to\_modal}).
\model{}'s physics, which is naturally pointwise (it depends on local $T$, $\qhum$,
$\pres_s$, and the underlying SST), is computed entirely on the nodal grid and
projected back to modal coefficients with \code{to\_modal} before being handed to
the integrator; the vorticity/divergence parts of a velocity tendency are formed
with the spectral operators \code{curl\_cos\_lat} and \code{div\_cos\_lat}
(\codref{chronos\_esm/atmos/dino\_atmos.py}{\code{\_forcing\_explicit}, return value}).

\begin{example}[Why a velocity seed must be made rotational first]
The symmetry-breaking perturbation (Section~\ref{sec:ch09eddies}) is built as a
random nodal velocity field $\uvec'$, multiplied by $\cos\lat$, transformed to
modal space, and then passed through \code{curl\_cos\_lat} to extract only its
rotational (vorticity) part:
\begin{lstlisting}
clu = jnp.stack([cl*amp*randn, cl*amp*randn])     # cos(lat)*(u',v')
seed_vort = grid.curl_cos_lat(grid.to_modal(clu)) # rotational part only
\end{lstlisting}
The seed is injected purely into \code{vorticity}; the divergent part is
discarded because divergent perturbations radiate away as gravity waves rather
than projecting onto the slowly growing baroclinic mode.
\end{example}

\section{Time stepping: implicit--explicit Runge--Kutta}

The primitive equations are \emph{stiff}: fast external and internal gravity waves
($\sim 300\ \mathrm{m\,s^{-1}}$) coexist with the slow advective flow
($\sim 10\ \mathrm{m\,s^{-1}}$). A fully explicit scheme would have to resolve the
gravity-wave CFL with a tiny step. The dycore instead splits the right-hand side
into the linear terms responsible for the fast waves (handled \emph{implicitly},
where they are unconditionally stable) and the nonlinear advective/forcing terms
(handled \emph{explicitly}). \model{} assembles this split as
\begin{equation}
  \dt{\mathbf{s}} \;=\; \underbrace{E(\mathbf{s})}_{\text{explicit}}
                      \;+\; \underbrace{I(\mathbf{s})}_{\text{implicit}},
  \qquad
  E(\mathbf{s}) = E_{\mathrm{dyn}}(\mathbf{s}) + F(\mathbf{s}),
\end{equation}
where $E_{\mathrm{dyn}}$ is \code{eq.explicit\_terms}, $F$ is the \model{} forcing
of Section~\ref{sec:ch09forcing}, $I$ is \code{eq.implicit\_terms}, and the
implicit solve uses \code{eq.implicit\_inverse}. The integrator is the
three-stage implicit--explicit Runge--Kutta scheme \code{imex\_rk\_sil3}
(``semi-implicit, low-storage, 3rd order''):
\begin{lstlisting}
ode  = ti.ImplicitExplicitODE.from_functions(
           explicit, eq.implicit_terms, eq.implicit_inverse)
step = ti.step_with_filters(ti.imex_rk_sil3(ode, dt), [self.diff])
\end{lstlisting}
(\codref{chronos\_esm/atmos/dino\_atmos.py}{\code{\_run\_interval}}). The time
step is $\Delta t = 20$~min, nondimensionalised through the dycore's unit system
(\code{specs.nondimensionalize}); $\approx 72$ steps make a model day
(\code{steps\_per\_day}). One ocean coupling interval calls \code{step(state, sst,
n\_days)}, which runs $72\,n_{\text{days}}$ inner steps with the SST held fixed and
jitted as a traced argument so the whole interval compiles once.

\begin{remark}[A unit-handling foot-gun]
All physics is done in plain \jax{} with scalar scale factors precomputed in
\code{\_\_init\_\_} (e.g. \code{t0\_seconds}, \code{v\_scale\_ms}); no \code{pint}
objects ever enter a jitted region. SI rates in $\mathrm{s^{-1}}$ are made
nondimensional by multiplying by \code{t0\_seconds}. When a saved state is
reloaded, \code{sim\_time} must be restored as \code{None}, not as the stored
scalar: the ImEx integrator forms tendencies with \code{sim\_time=None} and adds
them to the state by \code{tree\_map}, so a scalar clock would trigger an
\code{array + None} mismatch on the first resumed step
(\codref{chronos\_esm/atmos/dino\_atmos.py}{\code{load\_state}}).
\end{remark}

\section{Scale-selective hyperdiffusion}

The truncation discards scales smaller than the grid; energy that cascades toward
that cutoff must be removed or it piles up as grid-scale noise. \model{} applies a
$\nabla^4$ (del$^4$, $\text{order}=2$) hyperdiffusion as a post-step filter. Using
the Laplacian eigenvalue \eqref{eq:laplacian-eig}, mode $(l,m)$ is multiplied each
step by a damping factor that scales like
\begin{equation}
  \hat f_l^m \;\longleftarrow\; \hat f_l^m \,\Big[\,1 + c\,\big(\,l(l+1)\big)^{2}\,\Big]^{-1},
  \qquad
  c \;\propto\; \frac{\Delta t}{\tau\,[\,L(L+1)\,]^{2}} ,
\end{equation}
so the very top (grid-scale) mode is damped on the chosen $e$-folding time $\tau$
while the largest scales are left almost untouched
(\codref{chronos\_esm/atmos/dino\_atmos.py}{\code{ti.horizontal\_diffusion\_step\_filter}}).
The crucial tuning is $\tau$. \model{} sets $\tau = 6$~h, deliberately
\emph{weaker} than the $\tau\approx 2$~h common in idealized cores. At T31 the
baroclinic eddies live at fairly high total wavenumbers ($l\sim 10$--$15$); a
$2$~h diffusion damps \emph{those} modes on $1$--$15$ days --- competitive with
their $1$--$3$ day baroclinic growth --- which would suppress the eddies and with
them the eddy momentum flux that drives the surface westerlies. At $\tau = 6$~h
the eddies grow while grid-scale noise is still killed quickly.

\begin{remark}[The opposite lesson from the single-level model]
In the legacy single-level atmosphere
(Chapter~\ref{chap:ch10_atmosphere_singlelevel}) eddies were \emph{spurious} and
hyperdiffusion was kept strong to suppress them. Here eddies \emph{are} the
physics we want, so diffusion must be weak enough to let them grow. The two
regimes call for opposite settings; copying tuning from one to the other is a
trap.
\end{remark}

\section{Held--Suarez forcing with an SST-slaved equilibrium}\label{sec:ch09forcing}

\dino{} integrates the adiabatic, frictionless dynamics; the diabatic physics is
\model{}'s. The thermal forcing follows the Held--Suarez benchmark
\citep{heldsuarez1994}: a Newtonian relaxation of temperature toward a prescribed
radiative-equilibrium field $\Teq$, plus Rayleigh drag in the boundary layer. The
key \model{} departure is that $\Teq$ is \emph{anchored to the underlying sea
surface temperature}, so the atmosphere responds to the ocean.

\paragraph{Equilibrium temperature.} The classical Held--Suarez $\Teq$ uses a
fixed analytic equator--pole profile $315 - 60\sin^2\lat$. \model{} keeps the
benchmark's vertical lapse structure but replaces the surface anchor with the
local SST $T_s$:
\begin{equation}
  \Teq(\lat,\sig) \;=\; \Big(\tfrac{\sig\,\pres_s}{\pres_0}\Big)^{\!\kappa}
  \Big[\, T_s(\lat,\lon) \;-\; \Delta_z\,\ln\!\Big(\tfrac{\sig\,\pres_s}{\pres_0}\Big)\,\cos^2\lat \,\Big],
  \qquad
  \Teq \geq T_{\min},
  \label{eq:teq}
\end{equation}
with $\pres_0 = 10^5$~Pa, the static-stability parameter $\Delta_z = 10$~K, and a
floor $T_{\min} = 200$~K
(\codref{chronos\_esm/atmos/dino\_atmos.py}{\code{\_equilibrium\_temperature}}).
The $(\sig\pres_s/\pres_0)^\kappa$ factor converts the surface (potential)
anchor into actual temperature on each $\sig$ level; the $\Delta_z\ln(\cdot)$ term
sets the vertical lapse rate, tapered to the tropics by $\cos^2\lat$. Because $T_s$
is the coupled ocean's SST, the equator--pole \emph{temperature gradient that
drives the jet is whatever the ocean produces} --- with the honest consequence
that a too-weak modelled SST gradient yields weaker-than-observed surface
westerlies.

\paragraph{Newtonian relaxation and drag.} The temperature and momentum forcings
are
\begin{align}
  F_T \;&=\; -\,k_T(\lat,\sig)\,\big(\,T - \Teq\,\big), \label{eq:newt}\\
  F_{\uvec} \;&=\; -\,k_v(\sig)\,\uvec, \label{eq:drag}
\end{align}
with Held--Suarez vertical profiles built from
$\mathrm{cut}(\sig)=\max\!\big(0,\,(\sig-\sig_b)/(1-\sig_b)\big)$, $\sig_b=0.7$:
\begin{equation}
  k_T \;=\; k_a + (k_s - k_a)\,\mathrm{cut}(\sig)\,\cos^4\lat,
  \qquad
  k_v \;=\; k_f\,\mathrm{cut}(\sig),
\end{equation}
with $k_a = (40\ \mathrm{d})^{-1}$ (free-atmosphere radiative timescale),
$k_s = (4\ \mathrm{d})^{-1}$ (faster boundary-layer relaxation), and
$k_f = (1\ \mathrm{d})^{-1}$ (boundary-layer drag), all nondimensionalised
(\codref{chronos\_esm/atmos/dino\_atmos.py}{\code{\_kt}, \code{\_kv} profiles;
\code{\_forcing\_explicit}}). The drag acts only below $\sig_b$ (the boundary
layer) and the strong thermal relaxation $k_s$ acts there and equatorward, exactly
as in \citet{heldsuarez1994}.

\paragraph{Moisture and the ITCZ.} On top of dry Held--Suarez, \model{} carries a
prognostic specific-humidity tracer $\qhum$ with a surface evaporation source and a
large-scale condensation sink. Evaporation into the lowest layer uses a bulk
formula driven by SST and near-surface wind, capped to $1$--$25\ \mathrm{m\,s^{-1}}$:
\begin{equation}
  E \;=\; \rho_a\,C_E\,|\uvec_{\mathrm{sfc}}|\,\max\!\big(q_{\mathrm{sat}}(T_s) - \qhum_{\mathrm{bot}},\,0\big)
  \quad[\mathrm{kg\,m^{-2}s^{-1}}],
\end{equation}
and supersaturation is rained out on a $\tau_c = 3$~h timescale, latent-heating the
column:
\begin{equation}
  C \;=\; \frac{\max(\qhum - q_{\mathrm{sat}}(T),\,0)}{\tau_c},
  \qquad
  F_T \mathrel{+}= \frac{\Lv}{c_p}\,C,
  \qquad
  \pd{\qhum}{t} \mathrel{+}= E\,\frac{\grav}{\pres_s\,\Delta\sig} - C,
\end{equation}
with $q_{\mathrm{sat}}$ from Clausius--Clapeyron
(\codref{chronos\_esm/atmos/dino\_atmos.py}{\code{\_qsat}, \code{\_forcing\_explicit}
moisture block}). Column-integrated $C$ is reported as precipitation
(\codref{chronos\_esm/atmos/dino\_atmos.py}{\code{diagnostics}}). Moisture
converging in the tropics condenses and releases latent heat, producing an ITCZ
that the dry single-level model could never represent.

\begin{remark}[Conservation fixers]
Two pragmatic fixes keep long runs honest. (i) Surface pressure is rescaled every
interval to conserve the initial global-mean dry mass; without it, Gibbs ripples
from spectrally representing orography leak $\sim 0.3$~hPa/day
(\codref{chronos\_esm/atmos/dino\_atmos.py}{\code{\_fix\_mass}}). (ii) The ETOPO
orography is Gaussian-smoothed before use, because raw T31 topography has sharp
gradients that produced a $\sim 1.8$~hPa/day mass drift through those ripples.
Reported MSLP reduces $\pres_s$ to mean sea level with a fixed standard-atmosphere
scale height, so it is noise-free and comparable to ERA5.
\end{remark}

\section{How the jet and storm tracks emerge}\label{sec:ch09eddies}

The point of going multi-level is that the mid-latitude circulation is now
\emph{earned}, not prescribed. The SST-anchored $\Teq$ \eqref{eq:teq} sets up an
equator--pole temperature gradient; thermal-wind balance turns that gradient into
a westerly jet with vertical shear; the sheared, stratified flow is
baroclinically unstable; the growing eddies flux momentum into the jet core and
heat poleward, and their surface signature is the trade-easterly /
mid-latitude-westerly / polar-easterly tripole. None of this is imposed --- it is
the saturation of baroclinic instability, exactly the textbook life cycle
\citep[Ch.~9, 12]{vallis2017}.

Getting there in finite wall-clock time needed two non-obvious choices.

\paragraph{(1) Initialise near radiative equilibrium.} \dino{}'s
\code{isothermal\_rest\_atmosphere} is an isothermal fluid at rest --- and
\emph{exactly} zonally symmetric. \model{} instead initialises the temperature
field at the SST-anchored $\Teq$ so the baroclinic jet exists from day~0 rather
than taking $\sim 2$ months to build through the slow $40$-day thermal relaxation
(\codref{chronos\_esm/atmos/dino\_atmos.py}{\code{initial\_state}}).

\paragraph{(2) Break zonal symmetry with a \emph{rotational} seed.} An exactly
axisymmetric state is an unstable equilibrium the flow never leaves on its own ---
only floating-point roundoff breaks it, and only after $\sim 100$ days. Without an
asymmetric seed \emph{no} baroclinic eddies form, there is no eddy
momentum-flux convergence, and the surface stays easterly everywhere (the $u_{\mathrm{sfc}}$
pattern correlation against ERA5 is $\sim 0$). \model{} adds a small random
\emph{vorticity} perturbation ($\approx 2\ \mathrm{m\,s^{-1}}$ amplitude). It must
be rotational, not thermal: a temperature perturbation is simply relaxed away by
$F_T$ before it can grow, whereas a rotational seed projects directly onto the
growing baroclinic mode (\codref{chronos\_esm/atmos/dino\_atmos.py}{\code{\_\_init\_\_},
\code{seed\_vort} via \code{curl\_cos\_lat}}). With (1), (2), and the weak $6$~h
diffusion, eddies, storm tracks, and surface westerlies appear in $\sim 2$--$3$
weeks instead of $\sim 2$--$3$ months. The validated dry-core spin-up
(\codref{experiments/dino\_held\_suarez.py}{Phase-1 Held--Suarez verification})
reproduces the textbook twin $\sim 22\ \mathrm{m\,s^{-1}}$ upper-tropospheric
midlatitude jets with surface trades.

\begin{proposition}[The surface westerlies are eddy-driven]
In a statistically steady, zonally averaged state the surface zonal-momentum
budget below the boundary layer balances the eddy momentum-flux convergence
against surface drag,
\begin{equation}
  -\,\frac{1}{\Erad\cos^2\lat}\,\pd{}{\lat}\big(\overline{u'v'}\cos^2\lat\big)
  \;\approx\; k_v\,\overline{\uu}_{\mathrm{sfc}} .
\end{equation}
Hence $\overline{\uu}_{\mathrm{sfc}}>0$ (westerly) requires a convergence of eddy
momentum flux at that latitude. With the rotational seed switched off
(\code{seed\_wind\_ms=0}) the eddies vanish, $\overline{u'v'}\to 0$, and the
surface westerlies collapse to easterlies everywhere --- the failure mode the seed
was introduced to cure.
\end{proposition}

\section{Status and honest limitations}

The dino atmosphere is the validated control-run path through the standalone
harness \codref{experiments/run\_dino\_coupled.py}{checkpointing/resumable
control harness}, \emph{not yet} the path inside \code{main.py:step\_coupled}; the
single-level atmosphere of Chapter~\ref{chap:ch10_atmosphere_singlelevel} is still
the coupled default. Two honest caveats carry into the validation chapters
(\ref{chap:ch16_climatology}, \ref{chap:ch18_jets_stormtracks}). First, the
surface wind and pressure \emph{pattern} are only as good as the SST gradient the
T31 ocean supplies; a weak modelled gradient yields weaker westerlies than ERA5,
a limitation of the \emph{coupling}, not the dycore. Second, resume restores the
prognostic modal state exactly but recomputes diagnostic fields, so a restarted
run is state-complete but not bit-reproducible ($\sim 0.03$~K/week drift,
negligible). The ITCZ, by contrast, is a clear win over the single-level model:
moisture convergence makes it rain at the equator rather than in the subtropics.

\begin{exercise}[Switch off the symmetry-breaking seed]
Construct \code{DinoAtmosphere(seed\_wind\_ms=0.0)} and a second instance with the
default $2\ \mathrm{m\,s^{-1}}$ seed. Spin both up for $60$ days at a fixed
Earth-like SST and plot the zonal-mean surface zonal wind
$\overline{\uu}_{\mathrm{sfc}}(\lat)$ from \code{diagnostics}. Confirm that the
seed-free run has easterlies at all latitudes while the seeded run develops
mid-latitude westerlies. Relate the difference to the eddy momentum-flux budget of
the proposition above.
\end{exercise}

\begin{exercise}[The diffusion timescale controls the eddies]
Re-run a $60$-day spin-up at $\tau \in \{2,4,6,12\}$~h
(\code{diffusion\_tau\_hours}). For each, measure eddy kinetic energy
$\tfrac{1}{2}\,\overline{(u'^2+v'^2)}$ at $\sim 45^\circ$ and the peak surface
westerly. Show that $\tau=2$~h damps the eddies and weakens the westerlies, and
explain why this is the \emph{opposite} of the correct setting for the
single-level model.
\end{exercise}

\begin{exercise}[SST gradient sets the jet]
Replace the default SST $T_s = 300 - 40\sin^2\lat$ with a weaker gradient
$300 - 20\sin^2\lat$ and a stronger one $300 - 60\sin^2\lat$, passing each to
\code{initial\_state} and \code{step}. Using thermal-wind balance, predict how the
upper-tropospheric jet speed should change, then verify against the model's
$\overline{\uu}(\lat,\sig)$. Use this to argue why a cold-biased tropical SST in
the coupled run (a known \model{} bias) weakens the simulated westerlies.
\end{exercise}

\begin{exercise}[Differentiable sensitivity of the jet]
Because the whole trajectory is \jax{}-differentiable
(Chapter~\ref{chap:ch03_differentiable}), use \code{jax.grad} to compute the
sensitivity of a $30$-day-mean jet-speed diagnostic to the static-stability
parameter $\Delta_z$ (\code{dThz}) in \eqref{eq:teq}. Compare the adjoint gradient
against a finite-difference estimate and comment on what the sign of
$\partial(\text{jet})/\partial\Delta_z$ tells you about the static-stability
control of baroclinicity.
\end{exercise}
